\documentclass[12pt,a4paper]{article}
\usepackage{amsmath}
\usepackage{graphicx}
\usepackage{caption}
\usepackage{hyperref}
\usepackage{geometry}
\geometry{left=1in,right=1in,top=1in,bottom=1in}

\title{Vibration Analysis of a Free-Free Beam: A Detailed Mathematical Framework}
\author{Jacob Christopher Vaught}
\date{\today}

\begin{document}

\maketitle

\begin{abstract}
This report presents an in-depth analysis of the vibration characteristics of a free-free beam using Euler-Bernoulli beam theory. The analysis begins with fundamental assumptions and builds a mathematical framework to describe the system. The vibration equations are derived, and the \( k_n/\beta_n \) factors are computed. We then explore node and antinode locations, detailing the physical significance and derivation of these points.
\end{abstract}

\section{Introduction}
In this report, we explore the vibration behavior of a free-free beam under the assumptions of linear elasticity and small deformations. The beam parameters are:
\begin{itemize}
    \item Young's Modulus, \( E = 69 \times 10^9 \, \text{Pa} \)
    \item Density, \( \rho = 2700 \, \text{kg/m}^3 \)
    \item Width, \( b = 0.0127 \, \text{m} \)
    \item Height, \( h = 0.003175 \, \text{m} \)
    \item Length, \( L = 1 \, \text{m} \)
\end{itemize}

We will derive the fundamental vibration equations using these parameters and assumptions to predict the natural frequencies and mode shapes.

\section{Fundamental Assumptions}
To begin the analysis, we make the following assumptions:
\begin{itemize}
    \item \textbf{Linear Elasticity:} The material of the beam follows Hooke's Law, implying a linear relationship between stress and strain.
    \item \textbf{Small Deformations:} The beam undergoes small displacements, allowing us to linearize the equations of motion.
    \item \textbf{No Initial Motion:} At \( t = 0 \), the beam is at rest, implying zero initial velocity: 
    \[
    \left. \frac{\partial y}{\partial t} \right|_{t=0} = 0.
    \]
    \item \textbf{No Initial Displacement:} The initial displacement is zero across the length of the beam:
    \[
    y(x, 0) = 0 \quad \text{for} \quad 0 \leq x \leq L.
    \]
    \item \textbf{Free-Free Boundary Conditions:} Both ends of the beam are free to move and rotate, resulting in zero moment and shear force at \( x = 0 \) and \( x = L \).
\end{itemize}

\section{Derivation of the Equation of Motion}
The equation of motion for a vibrating beam is derived from the balance of forces and moments. The Euler-Bernoulli beam theory gives us the following partial differential equation for transverse vibrations:
\begin{equation}
    \frac{\partial^4 y(x,t)}{\partial x^4} + \frac{\rho A}{EI} \frac{\partial^2 y(x,t)}{\partial t^2} = 0.
\end{equation}
Here, \( y(x,t) \) is the transverse displacement, \( E \) is the Young's modulus, \( I = \frac{bh^3}{12} \) is the second moment of area, \( \rho \) is the density, and \( A = bh \) is the cross-sectional area.

### Separation of Variables
We assume a solution of the form:
\begin{equation}
    y(x,t) = w(x) T(t),
\end{equation}
where \( w(x) \) is the spatial component and \( T(t) \) is the temporal component.

Substituting into the equation of motion and separating variables, we obtain two ordinary differential equations:
\begin{align}
    &\frac{d^4 w(x)}{dx^4} = \beta^4 w(x), \quad \beta^4 = \frac{\rho A \omega^2}{EI}, \\
    &\frac{d^2 T(t)}{dt^2} + \omega^2 T(t) = 0.
\end{align}

The solution for \( T(t) \) is:
\begin{equation}
    T(t) = A \cos(\omega t) + B \sin(\omega t),
\end{equation}
where \( \omega \) is the natural frequency.

\section{General Solution for Spatial Part}
The general solution for the spatial part \( w(x) \) is:
\begin{equation}
    w(x) = C_1 \cos(\beta x) + C_2 \sin(\beta x) + C_3 \cosh(\beta x) + C_4 \sinh(\beta x).
\end{equation}

\subsection{Application of Boundary Conditions}
For a free-free beam, the boundary conditions are:
\begin{align}
    w''(0) = 0, & \quad w'''(0) = 0, \\
    w''(L) = 0, & \quad w'''(L) = 0.
\end{align}

These conditions imply that both the bending moment and shear force are zero at the beam ends.

### At \( x = 0 \):
\[
\begin{aligned}
    w''(0) &= -C_1 \beta^2 + C_3 \beta^2 = 0, \\
    w'''(0) &= -C_2 \beta^3 + C_4 \beta^3 = 0.
\end{aligned}
\]
This implies \( C_1 = C_3 \) and \( C_2 = C_4 \).

### At \( x = L \):
\[
\begin{aligned}
    w''(L) &= -C_1 \beta^2 \cos(\beta L) + C_2 \beta^2 \cosh(\beta L) = 0, \\
    w'''(L) &= C_1 \beta^3 \sin(\beta L) + C_2 \beta^3 \sinh(\beta L) = 0.
\end{aligned}
\]
These conditions lead to the characteristic equation:
\begin{equation}
    \cosh(\beta L) \cos(\beta L) = 1.
\end{equation}

The roots of this equation, \( \beta_n \), correspond to the natural frequencies.

\section{Determination of \( k_n/\beta_n \) Factors}
The wavenumber \( k_n \) is related to the natural frequencies as:
\begin{equation}
    k_n = \frac{\beta_n}{L}.
\end{equation}
Substituting into the frequency equation, we obtain:
\begin{equation}
    \omega_n = \beta_n^2 \sqrt{\frac{EI}{\rho A}}.
\end{equation}

\section{Mode Shapes and Node/Antinode Locations}
The mode shapes are described by the function:
\begin{equation}
    w_n(x) = \sin\left(\frac{\beta_n x}{L}\right) + \sinh\left(\frac{\beta_n x}{L}\right) + \alpha_n \left[ \cos\left(\frac{\beta_n x}{L}\right) + \cosh\left(\frac{\beta_n x}{L}\right) \right],
\end{equation}
where \( \alpha_n = \frac{-\sin(\beta_n) + \sinh(\beta_n)}{\cos(\beta_n) + \cosh(\beta_n)} \).

### Nodes and Antinodes
Nodes are the points where \( w_n(x) = 0 \), and antinodes are where \( w_n(x) \) reaches its maximum absolute value. They can be found by setting the derivative of \( w_n(x) \) with respect to \( x \) equal to zero and solving for \( x \).

\section{Discussion}
The assumptions made simplify the problem and allow us to obtain analytical solutions. Nodes and antinodes provide critical information about the vibration pattern of the beam. Understanding these points is essential for applications involving structural dynamics and vibration control.

\section{Conclusion}
We have derived the fundamental equations governing the vibration of a free-free beam from basic assumptions. The mode shapes and natural frequencies were calculated, and the locations of nodes and antinodes were identified. This detailed analysis offers a comprehensive understanding of the dynamic behavior of beams in various engineering applications.

\end{document}
